\documentclass[10pt,a4paper]{amsart}
\usepackage[margin=19mm]{geometry}
\usepackage{amsmath,amssymb,mathtools}
\usepackage[scr=rsfs,cal=euler]{mathalfa}
\usepackage{xfrac}
\usepackage[unicode,hidelinks,bookmarks=false]{hyperref}
\newcommand{\ie}{i.e.\ }
\newcommand{\cf}{cf.\ }
\newcommand{\resp}{respectively\ }
\newcommand{\Subsection}{\subsection}
\providecommand{\nobreakdash}{\nobreak\hbox{-}}
\setlength{\emergencystretch}{2em}
\multlinegap0pt
\usepackage[shortlabels]{enumitem}
\usepackage{xfrac}
\usepackage{nicefrac}
\usepackage{amssymb}
\usepackage{dsfont}
\usepackage{multirow}
\usepackage{xcolor}
\usepackage{booktabs}
\usepackage{mathtools}
\newtheorem{lemma}{Lemma}
\newcommand{\N}{\mathbb{N}}
\newcommand{\R}{\mathbb{R}}
\newcommand{\cA}{\mathcal{A}}
\newcommand{\from}{\colon} 
\newcommand{\norm}[2]{     \| #1       \|_{ #2 }}
\title{Regularity theory for a new class of fractional \\ parabolic stochastic evolution equations}
\author{Kristin Kirchner, Joshua Willems}
\date{Source: arXiv:2205.00248v2 (CC BY 4.0)}
\begin{document}
\maketitle

\noindent\emph{Source and licence.} Regularity theory for a new class of fractional \\ parabolic stochastic evolution equations. Version arXiv:2205.00248v2 (CC BY 4.0), distributed by arXiv under the Creative Commons Attribution 4.0 International licence, http://creativecommons.org/licenses/by/4.0/, as recorded in the arXiv licence metadata for this exact version and shown on the abstract page https://arxiv.org/abs/2205.00248v2. Full licence notice: \texttt{licenses/arxiv-2205.00248-CC-BY-4.0.txt}.\par\smallskip

\noindent\textit{Editorial scope.} Counted unit: the article's lemma lem:Bochner-space-operator (source0.tex lines 4499--4526). The article's complete original proof of it is delivered with it (source0.tex lines 4528--4601, one \texttt{proof} environment of 74 lines). No mathematics is written, repaired or re-derived: the statement and the proof are the article's own lines, in the article's own notation, and no retained line is substituted or re-flowed.  Retained uncounted as the source-local context the proof needs: lines 1088--1102.  Nothing was omitted from any retained span, so the package carries no omission record.  No retained line contains a citation, so the wrapper carries no bibliography and no bibliography entry.  No retained line contains a figure environment, an image-inclusion command or drawing code, so no image is delivered and the package declares no asset dependency.  Editorial formatting only, in addition: an AMS \texttt{amsart} preamble that restates the article's own declarations and macros used by the retained lines, namely the environments \texttt{lemma} on separate counters, together with the standard packages those lines need.  The retained environments print in this document as Lemma 1 in order of appearance, whereas the article prints the same items with the numbering of its own full document.  The complete native source of the article as distributed in the arXiv e-print for this exact version is bundled unchanged as \texttt{sources/125155/source0.tex}.
\begin{equation}\label{eq:def-cA} 
\begin{split}
&[\cA v](\vartheta)
:= Av(\vartheta), 
\quad\;\; 
v \in \mathsf D(\cA), 
\ 
\text{a.a.\ }\vartheta\in (0,T), 
\\
&\mathsf D(\cA) 
= L^2(0,T;\mathsf D(A)) 
:= 
\bigl\{ v \in L^2(0,T;H): \norm{ Av(\,\cdot\,)}{L^2(0,T;H)} < \infty \bigr\}.   
\end{split}
\end{equation}

\begin{lemma}\label{lem:Bochner-space-operator}
	Let $T\in(0,\infty)$ and 
	$A\from \mathsf D(A)\subseteq H\to H$ 
	be a linear operator on a real or complex 
	Hilbert space $H$.  
	Consider the associated operator 
	$\mathcal A$ on $L^2(0,T;H)$ 
	as defined in \eqref{eq:def-cA}.
	Then, the following hold: 
	%
	\begin{enumerate}[leftmargin=1cm, label={\normalfont(\alph*)}]
		\item\label{lem:Bochner-space-operator-a} 
		%
		$\mathcal A$ is bounded if and only if 
		$A$ is bounded, 
		and in that case we have 
		%
		\[
		\norm{\mathcal A}{\mathscr L(L^2(0,T;H))} 
		= 
		\norm{A}{\mathscr L(H)} ;
		\]  
		%
		\item\label{lem:Bochner-space-operator-b}
		%
		$\mathcal A$ is closed if and only if $A$ is.
	\end{enumerate}
\end{lemma}

\begin{proof}
	If $A$ is bounded, 
	then the inequality 
	$\norm{\mathcal A}{\mathscr L(L^2(0,T;H))} 
	\leq 
	\norm{A}{\mathscr L(H)}$ 
	is easily verified.
	Now suppose that $\mathcal A$ is bounded. 
	Then for all $x \in H$ we have
	%
	\begin{align*}
	\norm{Ax}{H} 
	&= 
	\norm{T^{-\nicefrac{1}{2}}\mathbf 1_{(0,T)} \otimes Ax}{L^2(0,T;H)} 
	= 
	\norm{\mathcal A(T^{-\nicefrac{1}{2}}\mathbf 1_{(0,T)} \otimes x)}{L^2(0,T;H)} 
	\\
	&\leq 
	\norm{\mathcal A}{\mathscr L(L^2(0,T;H))} 
	\norm{T^{-\nicefrac{1}{2}}\mathbf 1_{(0,T)} \otimes x}{L^2(0,T;H)} 
	= \norm{\mathcal A}{\mathscr L(L^2(0,T;H))} \norm{x}{H}.
	\end{align*} 
	%
	Here, given $f \from (0,T)\to \R$ and $x \in H$,
	the function $f \otimes x \from (0,T)\to H$ is 
	defined by $[f\otimes x](t) := f(t)x$ for all $t \in (0,T)$.
	We thus find that $A$ is bounded 
	with operator norm  
	$\norm{A}{\mathscr L(H)}
	\leq \norm{\mathcal A}{\mathscr L(L^2(0,T;H))}$,
	which finishes the proof of~\ref{lem:Bochner-space-operator-a}. 
	
	To prove 
	part~\ref{lem:Bochner-space-operator-b}, 
	first let $A$ be closed and let the sequence 
	$(v_n)_{n\in\N}$ in $ \mathsf D(\mathcal A)$ 
	be such that $v_n \to v$ and $\mathcal Av_n \to y$ 
	in $L^2(0,T;H)$. We need to prove that 
	$v \in \mathsf D(\mathcal A)$ and 
	$y = \mathcal Av$. Let $(v_{n_k})_{k\in\N}$ 
	be a subsequence such that 
	$v_{n_k} \to v$ 
	and $Av_{n_k} \to y$ in~$H$, 
	a.e.\ in $(0,T)$, 
	so that by the closedness of $A$ it follows that 
	$v(\vartheta) \in \mathsf D(A)$ and $y(\vartheta) = Av(\vartheta)$ 
	for a.a.\ $\vartheta\in(0,T)$. 
	From the latter we obtain that $y = \mathcal Av$, 
	which is meaningful since $v,y\in L^2(0,T;H)$ 
	yields that $v \in \mathsf D(\mathcal A)$.
	
	Now let $\mathcal A$ be closed and 
	let $(x_n)_{n\in \N}$ in $\mathsf D(A)$ 
	be such that $x_n \to x$ and $Ax_n \to y$ in $H$. 
	This implies the following convergences in 
	$L^2(0,T;H)$: 
	%
	\begin{align*}
	\mathbf 1_{(0,T)} \otimes x_n 
	\to \mathbf 1_{(0,T)}\otimes x, 
	\\ 
	\mathcal A(\mathbf 1_{(0,T)} \otimes x_n) 
	= 
	\mathbf 1_{(0,T)} \otimes Ax_n  
	\to \mathbf 1_{(0,T)} \otimes y. 
	\end{align*}
	% 
	Since $\mathcal A$ is closed, 
	we deduce that 
	$\mathbf 1_{(0,T)} \otimes x \in \mathsf D(\mathcal A)$ 
	and $\mathbf 1_{(0,T)} \otimes y = \mathcal A(\mathbf 1_{(0,T)} \otimes x)$, 
	from which we may conclude $x \in \mathsf D(A)$ and $y = Ax$. 
	Hence $A$ is closed.
\end{proof}

\end{document}
